\documentclass[10pt,titlepage]{article}
\usepackage[letterpaper,margin=0.75in]{geometry}

\usepackage{booktabs}
\usepackage{longtable}
\usepackage{tabularx}
\usepackage{enumitem}
\usepackage{amssymb}
\usepackage{array}
\usepackage{parskip}
\usepackage{caption}
\usepackage[destlabel]{hyperref}
\hypersetup{
    colorlinks,
    citecolor=green,
    filecolor=magenta,
    linkcolor=red,
    urlcolor=cyan
}
\setlength{\emergencystretch}{5em}


% ----- Numbered system tests (same list style as the SRS) -------------------
% Numbering is shared by every testlist.  Moving or inserting a test changes
% its displayed number and all \testref links automatically.
\makeatletter
\newlength{\reqlabelwidth}
\newlength{\reqlabelsep}
\setlength{\reqlabelsep}{0.75em}
\newsavebox{\reqlabelbox}
\newcommand{\reqsetwidth}[2]{\expandafter\gdef\csname reqwidth@#1\endcsname{#2}}
\newcommand{\reqlabel}[1]{%
  \sbox\reqlabelbox{\textbf{#1}}%
  \ifdim\wd\reqlabelbox>\csname reqmax@\reqkind\endcsname\relax
    \expandafter\xdef\csname reqmax@\reqkind\endcsname{\the\wd\reqlabelbox}%
  \fi
  \usebox\reqlabelbox}
\newenvironment{reqlist}[1]{%
  \def\reqkind{#1}%
  \@ifundefined{reqmax@#1}{\expandafter\gdef\csname reqmax@#1\endcsname{0pt}}{}%
  \@ifundefined{reqwidth@#1}%
    {\settowidth{\reqlabelwidth}{\textbf{ST-NFR-00:}}}%
    {\setlength{\reqlabelwidth}{\csname reqwidth@#1\endcsname}}%
  \begin{itemize}[leftmargin=\dimexpr\reqlabelwidth+\reqlabelsep\relax,
    labelwidth=\reqlabelwidth,labelsep=\reqlabelsep,
    align=left,itemsep=0.9em,parsep=0.35em]}%
  {\end{itemize}%
  \immediate\write\@auxout{\string\reqsetwidth{\reqkind}%
    {\csname reqmax@\reqkind\endcsname}}}
\makeatother

\newcounter{vnvtestnum}
\newenvironment{testlist}{\begin{reqlist}{STNFR}}{\end{reqlist}}
\newcommand{\testdef}[1]{%
  \refstepcounter{vnvtestnum}\label{test:#1}%
  \reqlabel{ST-NFR-\arabic{vnvtestnum}:}}
\newcommand{\testref}[1]{\hyperref[test:#1]{ST-NFR-\ref*{test:#1}}}

% As in the Hazard Analysis, published SRS links use the label as their named
% destination; display numbers are written explicitly because the documents
% are compiled separately. This URL matches the SRS link used by the HA.
\newcommand{\srsurl}{https://team-hindsight.github.io/ophthamology-game/SRS/SRS.pdf}
\newcommand{\srsnfrref}[2]{\href{\srsurl\#nameddest=#1}{NFR#2}}
\newcommand{\srsdoc}[1]{\href{\srsurl}{#1}}
\newcommand{\srstableref}[2]{\href{\srsurl\#nameddest=#1}{#2}}

% Local auxiliary-parameter entries are individually anchored: a reference
% jumps to the parameter's row, even in a long table spanning several pages.
\newcommand{\constdef}[1]{\hypertarget{const:#1}{\texttt{\detokenize{#1}}}}
\newcommand{\constref}[1]{\hyperlink{const:#1}{\texttt{\detokenize{#1}}}}

\input{../Comments.text}
\input{../Common.text}

\begin{document}

\title{System Verification and Validation Plan for \progname{}} 
\author{\authname}
\date{\today}
	
\maketitle

\pagenumbering{roman}

\begin{table}[!htbp]
\caption{Revision History} \label{TblRevisionHistory}
\begin{tabularx}{\textwidth}{llX}
\toprule
\textbf{Date} & \textbf{Developer(s)} & \textbf{Change}\\
\midrule
2026/10/30 & Mohammad Bilal & Completed section 3.3, individual reflection, \& contributed to sections 1, 4.2, 6.1, 6.2, 7 \\
\bottomrule
\end{tabularx}
\end{table}

\newpage

\tableofcontents

\newpage

\listoftables

\listoffigures
\wss{Remove this section if it isn't needed}

\newpage

\section{Symbols, Abbreviations, and Acronyms}

\renewcommand{\arraystretch}{1.2}
\begin{longtable}{@{}p{0.24\textwidth}p{0.69\textwidth}@{}}
\caption{Symbols, Abbreviations, and Acronyms}\label{tab:VnVAcronyms}\\
\toprule
\textbf{Symbol} & \textbf{Description}\\
\midrule
\endfirsthead
\multicolumn{2}{c}{\tablename\ \thetable{} -- continued from previous page}\\
\toprule
\textbf{Symbol} & \textbf{Description}\\
\midrule
\endhead
\midrule
\multicolumn{2}{r}{Continued on next page}\\
\endfoot
\bottomrule
\endlastfoot
AV & Audiovisual; synchronized narration and animation.\\
CPU & Central Processing Unit.\\
FPS & Frames per second.\\
LC & Likely Change, as defined in the SRS.\\
MG & Module Guide.\\
MIS & Module Interface Specification.\\
NFR & Non-Functional Requirement.\\
OS & Operating System.\\
SBOM & Software Bill of Materials.\\
SRS & Software Requirements Specification.\\
ST-NFR & System test for a non-functional requirement; numbered in order of appearance.\\
VnV & Verification and Validation.\\
\end{longtable}

For the remaining requirements terminology, see the definitions and
abbreviations table in the SRS.

\newpage

\pagenumbering{arabic}

This document ... \wss{provide an introductory blurb and roadmap of the
  Verification and Validation plan}

\section{General Information}

\subsection{Summary}

\wss{Say what software is being tested.  Give its name and a brief overview of
  its general functions.}

\subsection{Objectives}

\wss{State what is intended to be accomplished.  The objective will be around
  the qualities that are most important for your project.  You might have
  something like: ``build confidence in the software correctness,''
  ``demonstrate adequate usability.'' etc.  You won't list all of the qualities,
  just those that are most important.}

\wss{You should also list the objectives that are out of scope.  You don't have 
the resources to do everything, so what will you be leaving out.  For instance, 
if you are not going to verify the quality of usability, state this.  It is also 
worthwhile to justify why the objectives are left out.}

\wss{The objectives are important because they highlight that you are aware of 
limitations in your resources for verification and validation.  You can't do everything, 
so what are you going to prioritize?  As an example, if your system depends on an 
external library, you can explicitly state that you will assume that external library 
has already been verified by its implementation team.}

\subsection{Custom Documents}

\wss{Summarize the CDs (Custom Documents) that were tackled by this project.  CDs 
can include usability testing, code walkthroughs, user documentation, formal
proof, GenderMag personas, Design Thinking, etc.  CDs should have already
been approved by the course instructor as included in your problem statement.
You can use a pull request to update your CDs (in TeamComposition.csv or
Repos.csv) if your plan changes as a result of the VnV planning exercise.}

\subsection{Relevant Documentation}

\wss{Reference relevant documentation.  This will definitely include your SRS
  and your other project documents (design documents, like MG, MIS, etc).  You
  can include these even before they are written, since by the time the project
  is done, they will be written.  You can create BibTeX entries for your
  documents and within those entries include a hyperlink to the documents.}

\wss{Don't just list the other documents.  You should explain why they are relevant and 
how they relate to your VnV efforts.}

\newpage
\section{Plan}

\wss{Introduce this section.  You can provide a roadmap of the sections to
  come.}

\subsection{Verification and Validation Team}

\wss{Your teammates.  Maybe your supervisor.
  You should do more than list names.  You should say what each person's role is
  for the project's verification.  A table is a good way to summarize this information.}

\subsection{SRS Verification}

\wss{List any approaches you intend to use for SRS verification.  This may
  include ad hoc feedback from reviewers, like your classmates (like your
  primary reviewer), or you may plan for something more rigorous/systematic.}

\wss{If you have a supervisor for the project, you shouldn't just say they will
read over the SRS.  You should explain your structured approach to the review.
Will you have a meeting?  What will you present?  What questions will you ask?
Will you give them instructions for a task-based inspection?  Will you use your
issue tracker?}

\wss{Maybe create an SRS checklist?}

\subsection{Design Verification}\label{sec:DesignVerification}

Design verification will be carried out after the Module Interface Specification
(MIS) is completed. The purpose is to check that the architecture described in
the Module Guide (MG) and the detailed design in the MIS are consistent with
the requirements in the SRS and address the relevant concerns identified
in the Hazard Analysis. The design will be verified through three reviews:

\begin{enumerate}[leftmargin=*, itemsep=0.5em]
    \item \textbf{Structured Internal Review:} All members of the Hindsight team
    will review the MG and MIS, preferably in a group meeting. One team member
    will be appointed to lead the review and work through the checklist in
    Appendix~\ref{sec:DesignVerificationChecklist}. The team will discuss any
    concerns, record them as GitHub issues, and assign members to address them.

    \item \textbf{External Peer Review:} Members of at least one other capstone
    team will review the design as part of the course's milestone peer review
    process. The checklist in Appendix~\ref{sec:DesignVerificationChecklist}
    will be shared with the reviewers, but they will not be required to use it.
    Their findings will be submitted as GitHub issues for the team to review
    and address.

    \item \textbf{Teaching Assistant Feedback:} The teaching assistant (TA)
    will evaluate the design documents as part of the course deliverable
    assessment. Any feedback requiring changes will be recorded as GitHub
    issues and addressed by the team.
\end{enumerate}

For each review, GitHub issues will identify the relevant section of the design,
the problem found, and the change needed. The assigned team member will make
the correction, and another team member will check the update before the issue
is closed.

\subsection{Verification and Validation Plan Verification}

\wss{The verification and validation plan is an artifact that should also be
verified.  Techniques for this include review and mutation testing.}

\wss{The review will include reviews by your classmates}

\wss{Create a checklists?}

\subsection{Implementation Verification}

\wss{You should at least point to the tests listed in this document and the unit
  testing plan.}

\wss{In this section you would also give any details of any plans for static
  verification of the implementation.  Potential techniques include code
  walkthroughs, code inspection, static analyzers, etc.}

\wss{The final class presentation in CAS 741 could be used as a code
walkthrough.  There is also a possibility of using the final presentation (in
CAS741) for a partial usability survey.}

\subsection{Automated Testing and Verification Tools}

\wss{What tools are you using for automated testing.  Likely a unit testing
  framework and maybe a profiling tool, like ValGrind.  Other possible tools
  include a static analyzer, make, continuous integration tools, test coverage
  tools, etc.  Explain your plans for summarizing code coverage metrics.
  Linters are another important class of tools.  For the programming language
  you select, you should look at the available linters.  There may also be tools
  that verify that coding standards have been respected, like flake9 for
  Python.}

\wss{If you have already done this in the development plan, you can point to
that document.}

\wss{The details of this section will likely evolve as you get closer to the
  implementation.}

\subsection{Software Validation}

\wss{If there is any external data that can be used for validation, you should
  point to it here.  If there are no plans for validation, you should state that
  here.}

\wss{You might want to use review sessions with the stakeholder to check that
the requirements document captures the right requirements.  Maybe task based
inspection?}

\wss{For those capstone teams with an external supervisor, the Rev 0 demo should 
be used as an opportunity to validate the requirements.  You should plan on 
demonstrating your project to your supervisor shortly after the scheduled Rev 0 demo.  
The feedback from your supervisor will be very useful for improving your project.}

\wss{For teams without an external supervisor, user testing can serve the same purpose 
as a Rev 0 demo for the supervisor.}

\wss{This section might reference back to the SRS verification section.}

\newpage
\section{System Tests}

This section describes the system tests used to verify the functional and
non-functional requirements in the SRS. The functional tests check that
the main features of the game work as expected, while the non-functional tests
check qualities such as performance, offline use, maintainability, and security.
Each test explains what will be checked and how the result will be recorded.

\subsection{Tests for Functional Requirements}

\wss{Subsets of the tests may be in related, so this section is divided into
  different areas.  If there are no identifiable subsets for the tests, this
  level of document structure can be removed.}

\wss{Include a blurb here to explain why the subsections below
  cover the requirements.  References to the SRS would be good here.}

\subsubsection{Area of Testing1}

\wss{It would be nice to have a blurb here to explain why the subsections below
  cover the requirements.  References to the SRS would be good here.  If a section
  covers tests for input constraints, you should reference the data constraints
  table in the SRS.}
		
\paragraph{Title for Test}

\begin{enumerate}

\item{test-id1\\}

Control: Manual versus Automatic
					
Initial State: 
					
Input: 
					
Output: \wss{The expected result for the given inputs.  Output is not how you
are going to return the results of the test.  The output is the expected
result.}

Test Case Derivation: \wss{Justify the expected value given in the Output field}
					
How test will be performed: 
					
\item{test-id2\\}

Control: Manual versus Automatic
					
Initial State: 
					
Input: 
					
Output: \wss{The expected result for the given inputs}

Test Case Derivation: \wss{Justify the expected value given in the Output field}

How test will be performed: 

\end{enumerate}

\subsubsection{Area of Testing2}

...

\subsection{Tests for Nonfunctional Requirements}\label{sec:NFRSystemTests}

The tests in this section verify the non-functional requirements (NFRs)
defined in the SRS. These requirements describe how well the game
should operate, rather than the specific features it provides.
The tests are organized according to the NFR categories in the
SRS, with each test identifying the requirement it covers, the conditions
for testing, the expected result, and how the test will be performed. Test
parameters are listed in Section~\ref{sec:VnVParameters}.

For the Capstone course, the team will test on the limited physical devices
available and use Unity tools to simulate additional device configurations
where appropriate. The devices or simulations used will be recorded when
the tests are performed. Any additional changes to the testing scope for
the Capstone course will be explained in a \textit{Remark} under the
relevant test.

\subsubsection{Performance Requirements}\label{sec:PerformanceNFRTests}

\begin{testlist}

\item[\testdef{ST_Startup}]\textbf{Startup Time}

\textit{Type:} Non-functional, Dynamic, manual.

\textit{Requirements Covered:} \srsnfrref{NFR_Startup}{11}.

\textit{Initial State:} The latest build is installed on the test devices, and the app is closed.

\textit{Input/Condition:} Open the app from its icon at least \constref{STARTUP_TRIALS_MIN} times on each test device. Close the app fully between launches.

\textit{Output/Result:} The median time to reach the usable home screen is at most \constref{STARTUP_MAX_S} seconds on each device tested.

\textit{How test will be performed:} Record each launch and measure the time from tapping the icon to when the home screen is ready. Calculate the median for each device.

\item[\testdef{ST_Responsiveness}]\textbf{Touch Response}

\textit{Type:} Non-functional, Dynamic, manual with timing logs.

\textit{Requirements Covered:} \srsnfrref{NFR_Responsiveness}{12}.

\textit{Initial State:} The app is open on the test devices.

\textit{Input/Condition:} Make at least \constref{TOUCH_TRIALS_MIN} valid touches across the menus and games. Also test each type of incorrect answer and time the correct-answer reveal.

\textit{Output/Result:} On each device tested, at least \constref{RESPONSE_RATE_MIN}\% of touches show a response within \constref{RESPONSE_MAX_MS} milliseconds. Every correct-answer reveal takes at most \constref{ANSWER_REVEAL_MAX_S} second.

\textit{How test will be performed:} Use Unity timestamps to record each touch event and its first visible response. Calculate the percentage that meets the limit, and record the time taken to reveal the correct answer after an incorrect response.

\item[\testdef{ST_Smoothness}]\textbf{Animation Frame Rate}

\textit{Type:} Non-functional, Dynamic, automatic.

\textit{Requirements Covered:} \srsnfrref{NFR_Smoothness}{13}.

\textit{Initial State:} The latest build is installed on the test devices.

\textit{Input/Condition:} Play through each procedure walkthrough and practice game while recording frame times.

\textit{Output/Result:} On each device tested, the game runs at or above \constref{FRAME_RATE_MIN} frames per second for at least \constref{FRAME_RATE_COMPLIANCE_MIN}\% of the measured time.

\textit{How test will be performed:} Use the Unity Profiler to collect frame times throughout each walkthrough and game. Calculate the percentage of measured time that meets the required frame rate.

\item[\testdef{ST_AVSync}]\textbf{Narration and Animation Timing}

\textit{Type:} Non-functional, Dynamic, automatic with manual review.

\textit{Requirements Covered:} \srsnfrref{NFR_AVSync}{14}.

\textit{Initial State:} The latest narration and animations are included in the build.

\textit{Input/Condition:} Play each narrated step in every procedure walkthrough on the test devices.

\textit{Output/Result:} The narration starts within \constref{AV_SYNC_MAX_MS} milliseconds of its matching animation for every step.

\textit{How test will be performed:} Add temporary Unity logs at the start of each narration clip and its animation. Compare the timestamps and review a recording of any step that appears out of sync.

\item[\testdef{ST_AudioSafety}]\textbf{Audio Levels}

\textit{Type:} Non-functional, Static audio inspection.

\textit{Requirements Covered:} \srsnfrref{NFR_AudioSafety}{15}.

\textit{Review procedure:} One team member will export the final narration, sound effects, and music with the game's audio mixer settings applied, then compare their peak levels using Audacity or a similar audio editor. A second team member will check the measurements and mixer settings. Any sound above the narration peak will be adjusted and checked again.

\textit{Expected result and evidence:} No sound effect or music has a peak level above narration. The team will save the measurement results, mixer settings, and any issues corrected.

\item[\testdef{ST_Flashing}]\textbf{Flashing Content}

\textit{Type:} Non-functional, Dynamic, automated analysis with manual review.

\textit{Requirements Covered:} \srsnfrref{NFR_Flashing}{16}.

\textit{Initial State:} All animations, transitions, and reward effects are included in the build.

\textit{Input/Condition:} Record or export every animation and visual effect, including repeated or looping effects.

\textit{Output/Result:} No sequence contains more than \constref{FLASH_MAX_PER_S} flashes in any one-second period.

\textit{How test will be performed:} A team member will run the recordings through a photosensitive-epilepsy flash-analysis tool. A second team member will review any flagged sections. Save the results and rerun the analysis after fixing any failing effects.

\item[\testdef{ST_Reliability}]\textbf{Crash-Free Sessions}

\textit{Type:} Non-functional, Dynamic, manual.

\textit{Requirements Covered:} \srsnfrref{NFR_Reliability}{17}.

\textit{Initial State:} The current build is installed on the test devices and crash logs are available.

\textit{Input/Condition:} Complete at least \constref{TEST_SESSIONS_MIN} sessions across the test devices. Include normal gameplay, restarting levels, saving, closing and reopening, and suspend/resume.

\textit{Output/Result:} At least \constref{CRASH_FREE_RATE_MIN}\% of sessions finish without a crash or freeze.

\textit{How test will be performed:} Record the result of each session. Check crash logs and count the app as frozen if it does not respond to a normal action for \constref{FREEZE_OBSERVATION_S} seconds. Calculate the percentage of sessions completed without a crash or freeze.

\item[\testdef{ST_DeviceHeat}]\textbf{CPU Usage During Gameplay}

\textit{Type:} Non-functional, Dynamic, automatic.

\textit{Requirements Covered:} \srsnfrref{NFR_DeviceHeat}{18}.

\textit{Initial State:} The app is installed on the test devices.

\textit{Input/Condition:} Play continuously for at least \constref{HEAT_TEST_SESSION_MIN} minutes on each test device, moving between the different activities.

\textit{Output/Result:} Measured CPU utilization stays at or below \constref{CPU_UTILIZATION_MAX_PERCENT}\% throughout the session on each device tested.

\textit{How test will be performed:} Use the Unity Profiler and available device monitoring tools to record CPU utilization during gameplay. Record the highest measured value and any thermal warning.

\end{testlist}

\subsubsection{Operational and Environmental Requirements}\label{sec:OperationalNFRTests}

\begin{testlist}

\item[\testdef{ST_Offline}]\textbf{Offline Functionality}

\textit{Type:} Non-functional, Dynamic, manual.

\textit{Requirements Covered:} \srsnfrref{NFR_Offline}{19}.

\textit{Initial State:} The game is installed with all content, but has not been opened for the first time.

\textit{Input/Condition:} Turn on airplane mode and disable Wi-Fi. Run all functional tests without connecting to the internet.

\textit{Output/Result:} Every functional test passes without requiring a network connection.

\textit{How test will be performed:} Run the full set of functional tests on the test devices and record the results. Record any missing content, error message, or feature that cannot be used offline.

\item[\testdef{ST_Installability}]\textbf{Installation and First Launch}

\textit{Type:} Non-functional, Dynamic, manual.

\textit{Requirements Covered:} \srsnfrref{NFR_Installability}{20}.

\textit{Initial State:} The app is not installed on the test device.

\textit{Input/Condition:} Install the latest release build and open the app for the first time.

\textit{Output/Result:} The home screen can be reached without creating an account, granting permissions, or downloading more content.

\textit{How test will be performed:} Install the build on each test device, record the first launch, and note any prompt or extra step required before reaching the home screen.

\textit{Remark:} For this capstone, the build will be installed directly for testing rather than through a public app-store listing. The store-installation part of the SRS requirement will remain unverified unless the app is published.

\item[\testdef{ST_InstallSize}]\textbf{Installed Application Size}

\textit{Type:} Non-functional, Static build inspection.

\textit{Requirements Covered:} \srsnfrref{NFR_InstallSize}{21}.

\textit{Review procedure:} Once the iOS and Android release builds are ready, the team member responsible for each build will check the installed application size using the platform's storage information. A second team member will check the reported numbers against \constref{INSTALL_SIZE_MAX_MB} MB. The build version, measured size, and source of the measurement will be recorded.

\textit{Expected result and evidence:} The installed size of each build is no more than \constref{INSTALL_SIZE_MAX_MB} MB. The team will keep screenshots or build reports showing the size.

\textit{Remark:} The SRS fit criterion uses the size reported by each app store. For the capstone, the team will use locally measured installed size unless a store release is available, and record that difference in the results.

\item[\testdef{ST_ClinicalAccuracy}]\textbf{Clinical Review of Game Content}

\textit{Type:} Non-functional, Manual document review.

\textit{Requirements Covered:} \srsnfrref{NFR_ClinicalAccuracy}{22}.

\textit{Review procedure:} One team member will prepare the latest game walkthroughs, narration, and procedure guide, with the version of each item recorded. A paediatric ophthalmologist will review the content against the procedures used at the clinic. Another team member will record the feedback and any required changes. Updated content will be sent back for approval.

\textit{Expected result and evidence:} The ophthalmologist approves every level and procedure-guide entry for the version being released, with no unresolved concerns. The team will keep the content checklist and written approval.

\end{testlist}

\subsubsection{Maintainability and Support Requirements}\label{sec:MaintainabilityNFRTests}

\begin{testlist}

\item[\testdef{ST_Maintainability}]\textbf{Effort to Make a Likely Change}

\textit{Type:} Non-functional, Manual estimation review.

\textit{Requirements Covered:} \srsnfrref{NFR_Maintainability}{23}.

\textit{Review procedure:} Once the main game is implemented, each team member will estimate how many hours it would take to implement and test each likely change listed in the SRS, excluding clinical review. The team will discuss the estimates and agree on a final number for each change. These will be compared with the original development effort.

\textit{Expected result and evidence:} Each estimated change effort is less than \constref{MAINT_FRACTION} of the original development effort. The team will keep the individual estimates, agreed estimates, and comparisons.

\textit{Remark:} For the scope of the capstone, the team will estimate the effort rather than carry out every likely change. This gives an indication of maintainability but does not fully verify the SRS criterion based on recorded effort.

\item[\testdef{ST_Portability}]\textbf{iOS and Android Compatibility}

\textit{Type:} Non-functional, Dynamic, manual.

\textit{Requirements Covered:} \srsnfrref{NFR_Portability}{24}.

\textit{Initial State:} The latest build is ready and the test configurations have been recorded.

\textit{Input/Condition:} Run all functional tests on the supported test configurations and check the game's layouts and controls at screen sizes from \constref{SCREEN_MIN_IN} to \constref{SCREEN_MAX_IN} inches.

\textit{Output/Result:} All functional tests pass on the configurations tested, and the layouts and controls remain usable at the tested screen sizes.

\textit{How test will be performed:} Run the tests and record the device or configuration, OS version, screen size, and results. List any supported iOS or Android versions or screen sizes that could not be checked.

\item[\testdef{ST_DependencyTesting}]\textbf{Third-Party Dependency Checks}

\textit{Type:} Non-functional, Dynamic regression testing with dependency inspection.

\textit{Requirements Covered:} \srsnfrref{NFR_DependencyTesting}{25}.

\textit{Initial State:} The Unity project, dependency files, and release build are available.

\textit{Input/Condition:} Check the Unity version and all third-party package and plugin versions. Run all functional tests on the release build after each dependency change.

\textit{Output/Result:} Every dependency has an exact version, and all functional tests pass after each dependency change with no unresolved issue before release.

\textit{How test will be performed:} One team member will prepare the dependency list from the Unity Package Manager and package files. A second team member will check the versions against the project. After a version change, the team will run the functional tests and keep the results with the dependency list.

\end{testlist}

\subsubsection{Security Requirements}\label{sec:SecurityNFRTests}

\begin{testlist}

\item[\testdef{ST_DataProtection}]\textbf{Private Local Storage}

\textit{Type:} Non-functional, Dynamic, manual.

\textit{Requirements Covered:} \srsnfrref{NFR_DataProtection}{26}.

\textit{Initial State:} A test nickname, settings, and game progress have been saved on an unmodified test device.

\textit{Input/Condition:} Try accessing the saved files from the device file manager, a different app, and normal sharing options.

\textit{Output/Result:} The saved data cannot be read by another app or through device file-sharing features.

\textit{How test will be performed:} Confirm where Unity writes the saved data and inspect the app's storage and sharing settings. Record the access attempts and check that the game can still load its own saves.

\item[\testdef{ST_Integrity}]\textbf{Release Signing and Distribution}

\textit{Type:} Non-functional, Static release inspection.

\textit{Requirements Covered:} \srsnfrref{NFR_Integrity}{27}.

\textit{Review procedure:} If the app is prepared for public release, the team members responsible for the iOS and Android builds will check the signatures using Xcode and Android build tools. Another team member will review the signing results and confirm that the planned public distribution channels are the Apple App Store and Google Play Store. The team will save the signing reports and store-release records, if available.

\textit{Expected result and evidence:} Both release builds have valid signatures, and public distribution uses only the official stores.

\textit{Remark:} This test will likely be omitted during the capstone because publishing the app to the Apple App Store and Google Play Store is difficult within the course timeline. The team may check signing on test builds, but that would not verify the full release and distribution requirement.

\item[\testdef{ST_Privacy}]\textbf{Network Activity}

\textit{Type:} Non-functional, Dynamic, manual with network capture.

\textit{Requirements Covered:} \srsnfrref{NFR_Privacy}{28}.

\textit{Initial State:} The game is installed on a test device connected to a monitored Wi-Fi network.

\textit{Input/Condition:} Capture network traffic while launching the game, playing each level, changing settings, earning rewards, and closing and reopening the game.

\textit{Output/Result:} The game makes no outgoing network requests.

\textit{How test will be performed:} Capture traffic using a network monitor or test proxy and look for requests associated with the app. Also check the Unity project for enabled analytics or network services. Keep the capture and record any unexplained request.

\item[\testdef{ST_Minimization}]\textbf{Stored Data Fields}

\textit{Type:} Non-functional, Dynamic save-file inspection.

\textit{Requirements Covered:} \srsnfrref{NFR_Minimization}{29}.

\textit{Initial State:} The app is installed with a fresh save file.

\textit{Input/Condition:} Set a test nickname, complete activities, earn and spend coins, and change the available settings. Then inspect the stored files.

\textit{Output/Result:} The app stores only fields belonging to Local Preferences, Character Inventory, Game Progress, and Journey Progress in the SRS data dictionary.

\textit{How test will be performed:} One team member will read the saved data using a development build or Unity editor. A second team member will compare the fields against the \srstableref{tab:JourneyDataDictionary}{SRS data dictionary} and record any extra fields.

\item[\testdef{ST_Immunity}]\textbf{Dependency Vulnerability Check}

\textit{Type:} Non-functional, Static dependency inspection.

\textit{Requirements Covered:} \srsnfrref{NFR_Immunity}{30}.

\textit{Review procedure:} One team member will list the exact versions of Unity, packages, and plugins included in the release. The member will check these versions against known-vulnerability databases and supplier advisories. A second team member will review any high or critical findings and check that fixes or safe versions have been applied.

\textit{Expected result and evidence:} No included third-party library has a publicly known high or critical vulnerability at the time of release. The team will keep the dependency list, the dated scan or advisory results, and any related GitHub issues.

\end{testlist}

\subsection{Traceability Between Test Cases and Requirements}

\wss{Provide a table that shows which test cases are supporting which
  requirements.}

\newpage
\section{Unit Test Description}

This section will be completed once the MIS (detailed design document) has been finalized \cite{MIS}.

% \wss{This section should not be filled in until after the MIS (detailed design
%   document) has been completed.}

% \wss{Reference your MIS (detailed design document) and explain your overall
% philosophy for test case selection.}  

% \wss{To save space and time, it may be an option to provide less detail in this section.  
% For the unit tests you can potentially layout your testing strategy here.  That is, you 
% can explain how tests will be selected for each module.  For instance, your test building 
% approach could be test cases for each access program, including one test for normal behaviour 
% and as many tests as needed for edge cases.  Rather than create the details of the input 
% and output here, you could point to the unit testing code.  For this to work, you code 
% needs to be well-documented, with meaningful names for all of the tests.}

\subsection{Unit Testing Scope}

TBD

% \wss{What modules are outside of the scope.  If there are modules that are
%   developed by someone else, then you would say here if you aren't planning on
%   verifying them.  There may also be modules that are part of your software, but
%   have a lower priority for verification than others.  If this is the case,
%   explain your rationale for the ranking of module importance.}

\subsection{Tests for Functional Requirements}

TBD

% \wss{Most of the verification will be through automated unit testing.  If
%   appropriate specific modules can be verified by a non-testing based
%   technique.  That can also be documented in this section.}

% \subsubsection{Module 1}

% \wss{Include a blurb here to explain why the subsections below cover the module.
%   References to the MIS would be good.  You will want tests from a black box
%   perspective and from a white box perspective.  Explain to the reader how the
%   tests were selected.}

% \begin{enumerate}

% \item{test-id1\\}

% Type: \wss{Functional, Dynamic, Manual, Automatic, Static etc. Most will
%   be automatic}
					
% Initial State: 
					
% Input: 
					
% Output: \wss{The expected result for the given inputs}

% Test Case Derivation: \wss{Justify the expected value given in the Output field}

% How test will be performed: 
					
% \item{test-id2\\}

% Type: \wss{Functional, Dynamic, Manual, Automatic, Static etc. Most will
%   be automatic}
					
% Initial State: 
					
% Input: 
					
% Output: \wss{The expected result for the given inputs}

% Test Case Derivation: \wss{Justify the expected value given in the Output field}

% How test will be performed: 

% \item{...\\}
    
% \end{enumerate}

% \subsubsection{Module 2}

% ...

\subsection{Tests for Nonfunctional Requirements}

TBD

% \wss{If there is a module that needs to be independently assessed for
%   performance, those test cases can go here.  In some projects, planning for
%   nonfunctional tests of units will not be that relevant.}

% \wss{These tests may involve collecting performance data from previously
%   mentioned functional tests.}

% \subsubsection{Module ?}
		
% \begin{enumerate}

% \item{test-id1\\}

% Type: \wss{Functional, Dynamic, Manual, Automatic, Static etc. Most will
%   be automatic}
					
% Initial State: 
					
% Input/Condition: 
					
% Output/Result: 
					
% How test will be performed: 
					
% \item{test-id2\\}

% Type: Functional, Dynamic, Manual, Static etc.
					
% Initial State: 
					
% Input: 
					
% Output: 
					
% How test will be performed: 

% \end{enumerate}

% \subsubsection{Module ?}

% ...

\subsection{Traceability Between Test Cases and Modules}

TBD

% \wss{Provide evidence that all modules have been considered.}

\newpage

\section{Appendix}

\subsection{Symbolic Parameters}\label{sec:VnVParameters}

The first table lists the SRS constants used in these tests. Their values
come directly from the SRS. The second
table defines the few additional values used to run the tests.

\begin{longtable}{@{}p{0.35\textwidth}p{0.12\textwidth}p{0.44\textwidth}@{}}
\caption{SRS Constants}\label{tab:VnVSRSConstants}\\

\toprule
\textbf{Constant} & \textbf{Value} & \textbf{Description}\\
\midrule
\endfirsthead

\multicolumn{3}{c}{\tablename\ \thetable{} -- continued from previous page}\\
\toprule
\textbf{Constant} & \textbf{Value} & \textbf{Description}\\
\midrule
\endhead

\midrule
\multicolumn{3}{r}{Continued on next page}\\
\endfoot

\bottomrule
\endlastfoot
\constdef{ADULT_TEXT_MAX_GRADE} & 8 & Maximum reading grade for caregiver text\\
\constdef{ANSWER_REVEAL_MAX_S} & 1 & Maximum answer reveal time (s)\\
\constdef{AV_SYNC_MAX_MS} & 100 & Maximum audiovisual start offset (ms)\\
\constdef{CLINICIAN_FIND_MAX_MIN} & 1 & Maximum time to find content for a patient (min)\\
\constdef{CLINICIAN_TRAINING_MAX_MIN} & 5 & Maximum clinician training time (min)\\
\constdef{CONTROL_POSITION_TOLERANCE_MM} & 5 & Maximum control position deviation (mm)\\
\constdef{CORE_TASK_MAX_MIN} & 5 & Maximum core task duration (min)\\
\constdef{CPU_UTILIZATION_MAX_PERCENT} & 80\% & Maximum CPU utilization\\
\constdef{CRASH_FREE_RATE_MIN} & 99\% & Minimum crash-free session rate\\
\constdef{FLASH_MAX_PER_S} & 3 & Maximum flashes per second\\
\constdef{FRAME_RATE_COMPLIANCE_MIN} & 95\% & Minimum time meeting frame-rate target\\
\constdef{FRAME_RATE_MIN} & 30 & Minimum frames per second\\
\constdef{HEAT_TEST_SESSION_MIN} & 30 & Minimum device heat-test duration (min)\\
\constdef{INSTALL_SIZE_MAX_MB} & 300 & Maximum installed application size (MB)\\
\constdef{LOW_SPEC_AGE_YEARS} & 5 & Minimum age of low-spec device (years)\\
\constdef{MAINT_FRACTION} & 0.10 & Maximum change effort / initial development effort\\
\constdef{MIN_ANDROID} & 10 & Minimum supported Android version\\
\constdef{MIN_IOS} & 16 & Minimum supported iOS version\\
\constdef{RESPONSE_MAX_MS} & 100 & Maximum time to respond to an input (ms)\\
\constdef{RESPONSE_RATE_MIN} & 95\% & Minimum quick-response rate\\
\constdef{SCREEN_MAX_IN} & 13 & Maximum supported screen diagonal (in)\\
\constdef{SCREEN_MIN_IN} & 4.7 & Minimum supported screen diagonal (in)\\
\constdef{SENTENCE_MAX_WORDS} & 10 & Maximum words per narrated sentence\\
\constdef{SKIP_HOLD_MIN_S} & 1 & Minimum press-and-hold time to skip (s)\\
\constdef{STARTUP_MAX_S} & 5 & Maximum median startup time (s)\\
\constdef{STARTUP_TRIALS_MIN} & 10 & Minimum cold-start trials per device\\
\constdef{TEST_SESSIONS_MIN} & 100 & Minimum test sessions for reliability\\
\constdef{TOUCH_TARGET_MIN} & 15 & Minimum touch target size (mm)\\
\constdef{UNAIDED_RATE_MIN} & 80\% & Minimum unaided task completion rate\\
\end{longtable}

\begin{longtable}{@{}p{0.49\textwidth}p{0.12\textwidth}p{0.30\textwidth}@{}}
\caption{Additional System-Test Parameters}\label{tab:VnVAdditionalParameters}\\

\toprule
\textbf{Constant} & \textbf{Value} & \textbf{Description}\\
\midrule
\endfirsthead

\multicolumn{3}{c}{\tablename\ \thetable{} -- continued from previous page}\\
\toprule
\textbf{Constant} & \textbf{Value} & \textbf{Description}\\
\midrule
\endhead

\midrule
\multicolumn{3}{r}{Continued on next page}\\
\endfoot

\bottomrule
\endlastfoot
\constdef{TOUCH_TRIALS_MIN} & 40 & Valid touch samples per device\\
\constdef{FREEZE_OBSERVATION_S} & 10 & Time allowed before recording a freeze (s)\\
\end{longtable}

\subsection{Design Verification Checklist}\label{sec:DesignVerificationChecklist}

This checklist will be used during the structured internal review described in
Section~\ref{sec:DesignVerification}. Reviewers will compare the MG and MIS
against the SRS and Hazard Analysis. The checklist is subject to change once the
MG and MIS are finalized.

\begin{itemize}[label={$\square$}, leftmargin=*, itemsep=0.35em]
    \item Are the system's main modules identified in the MG, with a clear
    purpose and responsibility for each one?

    \item Does the MIS provide enough detail about each module's inputs,
    outputs, interfaces, and expected behaviour for another team member to
    implement it?

    \item Are the modules and their dependencies consistent between the MG
    and MIS, without conflicting or unexplained relationships?

    \item Can the functional and non-functional requirements in the SRS be
    traced to the modules or design decisions that support them?

    \item Does the design cover the main gameplay flows, including selecting
    a level, completing it, saving progress, and exiting early?

    \item Does the design account for the safety-related behaviour in the
    Hazard Analysis, such as avoiding punitive feedback, preventing accidental
    input, and stopping audio and animations when a level is exited?

    \item Is the design divided into modules that can be tested and updated
    without unnecessary changes to unrelated parts of the game?

    \item Are external packages and assets identified so their licensing and
    any relevant security or content concerns can be checked?

    \item Are diagrams, interface descriptions, and assumptions clear and
    consistent enough for the team to implement and test the design?
\end{itemize}

\subsection{Usability Survey Questions?}

\wss{This is a section that would be appropriate for some projects.}

\newpage{}

\section{Appendix --- Generative AI Use Disclosure}

\subsection{AI Tools Used}

\begin{itemize}
    \item ChatGPT (GPT-6.1 Sol)
\end{itemize}

\subsection{How and Where Used}

\begin{itemize}
  \item In Section~\ref{sec:DesignVerification}, the team developed the design
  verification approach. GPT-6.1 Sol was used to improve the wording, grammar,
  and flow while preserving the team's original ideas and decisions.

  \item In Sections~\ref{sec:PerformanceNFRTests}--\ref{sec:SecurityNFRTests},
  GPT-6.1 Sol was used to draft the non-functional system tests. The team then
  reviewed and revised the test procedures, expected results, and scope remarks
  to align with the SRS and the resources available for the capstone.

  \item In Section~\ref{sec:VnVParameters}, GPT-6.1 Sol was used to help
  organize the symbolic parameter tables and link constants to the tests.
  The team verified the values against the SRS.

  \item In Section~\ref{sec:DesignVerificationChecklist}, GPT-6.1 Sol was used
  to draft the design verification checklist, which the team then reviewed
  and revised.

  \item GPT-6.1 Sol was also used to assist with automatic test numbering,
  cross-document references, and general \LaTeX{} formatting.
\end{itemize}

\subsection{Verification of AI-Generated Content}

The team reviewed the AI-assisted content against the SRS, Hazard Analysis,
design documents, and the scope of the project. The proposed verification and
validation activities were checked to ensure that the selected methods were
appropriate, the roles and review procedures were clear, and the planned
evidence would be sufficient to assess the requirements. The system and unit
test plans were checked for requirement coverage, traceability, clearly defined
inputs and expected results, and feasibility within the time and equipment
available to the team. Where a test could not fully meet the SRS fit criterion
within the capstone scope, the limitation was stated explicitly.

The team also checked AI-suggested thresholds, symbolic parameters, and
requirement references against the source documents. Checklists and review
procedures were examined for clarity and completeness. Finally, the \LaTeX{}
document was reviewed for correct numbering, cross-references, tables,
citations, and formatting. The team remains responsible for the accuracy,
completeness, consistency, and appropriateness of the final VnV Plan.

\newpage{}

\section{Appendix --- Reflection}

\wss{This section is not required for CAS 741}

The information in this section will be used to evaluate the team members on the
graduate attribute of Lifelong Learning.

\input{../Reflection.text}

\begin{enumerate}
  \item What went well while writing this deliverable? 
  \item What pain points did you experience during this deliverable, and how
    did you resolve them?
  \item What knowledge and skills will the team collectively need to acquire to
  successfully complete the verification and validation of your project?
  Examples of possible knowledge and skills include dynamic testing knowledge,
  static testing knowledge, specific tool usage, Valgrind etc.  You should look to
  identify at least one item for each team member.
  \item For each of the knowledge areas and skills identified in the previous
  question, what are at least two approaches to acquiring the knowledge or
  mastering the skill?  Of the identified approaches, which will each team
  member pursue, and why did they make this choice?
\end{enumerate}

\newpage
\bibliographystyle{unsrt}

\bibliography{../../refs/References}

\end{document}